\appendix


\section{\texorpdfstring{Proof of~\eqref{decreasingT}}{Proof of (6)}} 
\eqref{decreasingT}: If $s_i(\mu)<\mu$, and $u=(i,\mu_{i}+1)$, then
\[ T_iE_\mu = \frac{q^{\ell(u)+1}t^{a(u)}(1-t)}{1-q^{\ell(u)+1}t^{a(u)}}E_{\mu} - \frac{(t-q^{\ell(u)+1}t^{a(u)})(1-q^{\ell(u)+1}t^{a(u)+1})}{(1-q^{\ell(u)+1}t^{a(u)})^2} E_{s_i(\mu)},\]
where all arms and legs are in terms of the diagram of $s_i(\mu)$.
\begin{proof}
Suppose $\mu_i < \mu_{i+1}$, so  $s_i(\mu)<\mu$. Using~\eqref{quadratic}, note that $T_i^{-1} = \frac{1}{t}(T_i+1-t)$. Then since $s_i(\mu)<s_i(s_i(\mu))$, we can apply~\eqref{E:T-on-E-up} with $s_i(\mu)$ to obtain,

\[ T_i E_{s_i(\mu)} = E_\mu + \frac{t-1}{1-q^{\ell(u)+1}t^{a(u)}} E_{s_i(\mu)},\]
where the arms and legs are with respect to the diagram of $s_i(\mu)$. Multiplying both sides on the left by $T_i^{-1}$ and simplifying, we get,
\begin{align*}
E_{s_i(\mu)} &= \frac{1}{t}(T_i+1-t)(E_\mu + \frac{t-1}{1-q^{\ell(u)+1}t^{a(u)}}E_{s_i(\mu)})\\
tE_{s_i(\mu)} &= T_i E_\mu + (1-t) E_\mu + \frac{t-1}{1-q^{\ell(u)+1}t^{a(u)}} \left(T_i E_{s_i(\mu)} + (1-t) E_{s_i(\mu)} \right)\\
tE_{s_i(\mu)} &= T_i E_\mu + (1-t) E_\mu \\ &\hspace*{1.5cm} +\frac{t-1}{1-q^{\ell(u)+1}t^{a(u)}} \left( E_\mu + \frac{t-1}{1-q^{\ell(u)+1}t^{a(u)}} E_{s_i(\mu)} + (1-t) E_{s_i(\mu)} \right)\\
T_i E_\mu &= \left(t-1+\frac{1-t}{1-q^{\ell(u)+1}t^{a(u)}} \right) E_\mu \\ &\hspace*{1.5cm} + \left(t-\frac{(t-1)^2}{(1-q^{\ell(u)+1}t^{a(u)})^2} - \frac{(t-1)(1-t)}{1-q^{\ell(u)+1}t^{a(u)}} \right) E_{s_i(\mu)}\\
T_iE_\mu &= \frac{q^{\ell(u)+1}t^{a(u)}(1-t)}{1-q^{\ell(u)+1}t^{a(u)}}E_{\mu} - \frac{(t-q^{\ell(u)+1}t^{a(u)})(1-q^{\ell(u)+1}t^{a(u)+1})}{(1-q^{\ell(u)+1}t^{a(u)})^2} E_{s_i(\mu)}\qedhere
\end{align*}
\end{proof}


\section{Proof of Lemma~\ref{nonzero}}

\textbf{Lemma~\ref{nonzero}:} If $I_1$ is maximal with respect to $\gammalambdatilde$, then $\pIoneevaluated \neq 0$.

\begin{proof}
We will prove the contrapositive, so assume $p_{I_1}\left( \overline{\etamutilde}\right) = 0$. This implies either \begin{enumerate} \item $\overline{\eta}_{t_u} = t \overline{\eta}_j \text{ for some } u\in [1,r] \text{ and } j\in [t_u+1,t_{u+1}-1]$, \text{ or } \item $ q^{-1} \overline{\mutilde}_{m+1} = t \overline{\eta}_j \text{ for some } j \in [1,t_1-1]$.
 \end{enumerate}
 In the first case, if $\overline{\eta}_{t_u} = t\overline{\eta}_j$, the powers of $q$ being equal means columns $t_u$ and $j$ have the same height, so $\eta_{t_u} = \eta_j$. Let $I_1' := I_1 \cup \{j\}$, and let $\dI'$ be the action associated with $I_1'$. We will show that $d_I\gammanu = d_I'\gammanu$. In that context, we have $\gamma_j = \gamma_{t_{u+1}}$. The only positions affected by the addition of $j$ to $I_1$ are the positions $j$ and $t_1$. And in these positions,
 \[ (\dI \gammanu)_{t_u} = \gammanu_{t_{u+1}} \, \text{ and }\,  (\dI' \gammanu)_{t_u} = \gammanu_j\]
 \[ (\dI \gammanu)_{j} = \gammanu_j  \, \text{ and } \, (\dI' \gammanu)_{j} = \gammanu_{t_{u+1}}.\]
 All entries of $d_I\gammanu$ and $d_I'\gammanu$ agree, so $I_1$ is not maximal.
 
 In the second case, since $ q^{-1} \overline{\mutilde}_{m+1} = t \overline{\eta}_j \text{ for some } j \in [1,t_1-1]$, again using column height, we have $\eta_j = \mutilde_{m+1}-1$. In terms of $\gammanu$, because $\mutilde_{m+1} = \mutilde_{v_s+1}$, this implies $\gamma_j = \gamma_{t_1}$ after applying the definition of $\dI$. By the same argument as in the first case, $d_I\gammanu$ and $d_I'\gammanu$ are identical, hence $I_1$ is not maximal with respect to~$\gammanu$.
 \end{proof}

